\scanpage{16}

% 1–3, definition
Un ens ordonné \inserted{$D$} est appelé ``géométrie de drapeaux'' si pour tout $d\in D$, l'ens. \deleted{\unsure{dessous|dessiné|ordonné}} \inserted{\unsure{induit|initial}} $D(d)=\{d'\in D\mid d'\leq d\}$ a la structure d'un ``simplexe ordonné'' i.e. est isomorphe à $\unsure{\mathscr P^*|\mathcal P^*|\mathfrak P^*}(I_d)$,

% 2–3, upper right, directed to the power set
\margin{parties finies $\ne\varnothing$}

% 3–4, outer left
% slanted
\margin{\uline{NB} on aura donc $I_d\ne\varnothing$}

% 4–8, minimal elements and length
avec l'inclusion --- où $I_d$ peut s'identifier à l'ens des él. minimaux de $D(d)$, ou encore : la partie de l'ens $\Phi\unsure{=|.}D(\underline{1})$ des él. minimaux de $D$ (``drapeaux de longueur 1 --- NB $\operatorname{card}(I_d)$ s'appellera ``longueur'' du drapeau $d$).

% 8–11, antifacteur
Pour toute partie finie $\sigma$ de $\Phi$, soit $D(\sigma)$ l'ens des $d\in D$ tels que $\deleted{\unsure{\sigma\subset}}I_d\deleted{\unsure{\ni}}=\sigma$, on trouve un antifoncteur
\[
 (\text{parties finies \inserted{$\ne\varnothing$} de }\Phi)^o\xrightarrow{D}(\mathrm{Ens})\qquad
 \text{tel que }\operatorname{card}D(\{i\})=1\quad\forall i\in\Phi
\]

% 12–14, reconstruction
dont la connaissance permet de reconstituer $D_{\unsure{*|{}}}$ comme
\[
 D_{\unsure{*|{}}}\simeq\coprod_{\sigma\in\mathscr P_f^*(\Phi)}D(\sigma)
\]
avec la relation d'ordre évidente. Ainsi

% 15–18, equivalence
\[
\begin{aligned}
\left\{\begin{array}{l}\text{géométrie de drapeaux}\\\text{(ens. ordonné \unsure{abs.|avec|ainsi})}\end{array}\right.
 &\Longleftrightarrow \begin{array}{l}
 \text{Couple }(\Phi,D)\text{ d'un ens}\\
 \Phi\text{ et d'un foncteur}\\
 \underline{\mathscr P_f^*}(\Phi)^o\mathbin{-}(\mathrm{Ens})\\
 \text{tel que }\operatorname{card}D(\{i\})=\unsure{4|1}\quad\forall i\in\unsure{I|\Phi}
 \end{array}
\end{aligned}
\]

% 17–19, lower left of equivalence
\margin{\uline{NB} $\Phi$ est tjrs l'ens des sommets d'une maquette simpliciale, mais bien sûr $D$ n'est pas déterminé par celle-ci\ldots}

% 19–20, strictness
\uline{NB} Pour que dans $D$ on ait $\forall d\in D$
\[
 d=\mathop{\rm Sup}_{i\in I_d}i
\]

% 19–20, right of supremum formula
\margin{géom de drapeaux stricte}

% 21–24, subsets
il f. et s. que $D$ soit isomorphe à un ens des parties finies \inserted{$\ne\varnothing$} de $\Phi$ avec l'inclusion (contenant les parties finies ponctuelles et contenant, avec tout ens, les ens $\ne\varnothing$ plus petits). On trouve ainsi

% 25, correspondence
\[
 \left(\begin{array}{c}\text{géom. de drapeaux}\\\text{stricte}\end{array}\right)
 \Longleftrightarrow
 \left(\begin{array}{c}\text{maquettes}\\\text{simpliciales}\end{array}\right)
\]

% 26–28, functor formulation
En termes du foncteur $D:\mathscr P_f^{\unsure{\triangle|*}}(\Phi)^o\mathbin{-}(\mathrm{Ens})$, celà signifie que $D$ est \deleted{le} \inserted{\unsure{sous|isom. à un|un sous-}} foncteur \deleted{\unsure{prenant}} \inserted{du foncteur} \inserted{final i.e.} \deleted{qui \unsure{en chaque} $\sigma$, \unsure{a pour}} $\forall\sigma,$ $\operatorname{card}D(\sigma)\leq1$ i.e. $\operatorname{card}D(\sigma)\in\{0,1\}$

% lower left
% archive sheet number: 15
